\section{Related Work}
%\ka{We should also mention the weakness of using very narrow camera position. For example, in some applications like UAV or underwater imaging it is desirable to be able to take images from a wider area. Also, some points on why a spread-out camera position is better than fixed elevation images despite a decrease in reconstruction performance}
%\sk{Certainly, we can maybe push this to intro to highlight the gaps in literature? Can you write up a gist for this along with a reason for why specifically in underwater and UAV imaging?} 
%\ka{I will add this in the dataset section}
% \subsection{3D Reconstruction}

% Classical methods

% Deep learning based
Early works on 3D reconstruction focus on mathematically modeling the imaging process and solving an ill posed problem via suitable scene priors and geometric constraints \cite{hartley2003multiple}.
Among the recent deep learning based methods, there are three broad types of techniques based on the choice of output representation :
Volume-based, Surface-based and Intermediate representation.
Surface based methods explore meshes and point clouds that are unordered representations \cite{wang2018pixel2mesh, kanazawa2018learning, pontes2018image2mesh}. 
Volumetric methods explore regular voxel grid based representation as we consider in this paper \cite{wu2017marrnet, tulsiani2017multi, yang2018learning}.
The recent state of the art works in this category include Pix3D \cite{sun2018pix3d}, 3DR2N2 \cite{3dr2n2} and Pix2Vox \cite{pix2vox} which we use as our baselines to compare with.
There are also techniques that decompose the problem into sequential steps each predicting a different representation that can be used as per application \cite{tatarchenko2016multi, lin2018learning}.

% \subsection{Next View Selection}

% Use active cameras
Early works in NVS acquire range images and aim to recover surface geometry of an object.
Connolly et al. \cite{connolly1985determination} did one of the first works on automatic range sensor positioning using partial octree models and optimize for surface area covered.
Maver et al. employ occlusion guidance \cite{maver1993occlusions} and optimize for information gained \cite{maver1993planning} for NVS.
These and related methods \cite{pito1996sensor, pito1995solution, pito1999solution} require the images to have an overlap to allow for registration and integration of the new data captured with previous scans. 
Earliest work in data driven NVS is by Banta et al. \cite{banta1995best} which use synthetic range data to predict occupancy grids.
Here, the number of viewpoints on the view sphere are limited and the radius is large enough to encompass the whole object.
The center and coordinate axes of the view sphere are aligned with those of the object and the camera always faces this centre.
We work with a similar setup but use passive RGB cameras.

Recent works also focus on active sensors such as Hepp et al. \cite{hepp2018learn} use a CNN to learn a viewpoint utility function, Delmarico et al. \cite{delmerico2018comparison} optimize volumetric infromation gain and Mendez et al. \cite{mendez2017taking} maximize information gain in terms of total map coverage and reconstruction quality.
Chen et al. \cite{chen2005vision} propose a prediction model for the unknown area of the object using a 3D range data acquisition.
Zhang et al. \cite{delmerico2018comparison} use self oclusion information for NVS to predict 3D motion estimation.
Potthast et al. \cite{potthast2014probabilistic} and Isler et al. \cite{isler2016information} propose information gain based NVS for occluded enviroments and 3D reconstruction respectively.
Other works include greedy \cite{obwald2018gpu} and supervised learning approach \cite{ly2019autonomous} for visibility based exploration and reconstruction 3D enviroments.
Hashim et al. \cite{hashim2019next} propose a view selection recommendation method by minimizing pose ambiguity for object pose estimation using depth sensors.
There are also works on robotic arm navigation for autonomous exploration and grasping \cite{monica2018contour, morrison2019multi, bai2017toward, wu2019plant}, UAV positioning in remote sensing \cite{almadhoun2019guided}, Pose estimation \cite{sock2017multi} and 3D scene inpainting \cite{han2019deep}.

\begin{figure*}
\begin{center}
\includegraphics[scale=0.43]{figures/cvpr_architecture_final.pdf}
\end{center}
   \caption{The proposed 3D-NVS architecture showing various constituent stages: NVS-Net, Encoder, Decoder and Refiner modules. Note that Encoder shares weight with NVS-Net in the 3D-NVS model.}
\label{fig:model-architecture}
\end{figure*}

Closely related to our work and among the deep learning based methods, Wu et all. \cite{wu20153d} propose a deep neural network based view prediction but use 2.5D RGB-D images to improve 3D object recognition.
Daudelin et al. \cite{daudelin2017adaptable} propose a probabilistic NVS algorithm but use RGB-D sensor data for 3D object reconstruction.
Mendez et al. \cite{mendez2016next} use stereo image pairs and propose NVS to optimize for surface area covered by reconstructing object surface.
Dunn et al. \cite{dunn2009next} use passive cameras but optimize for surface area covered as their goal is surface reconstruction.
Mendoza et al. \cite{mendoza2020supervised} also pose the this as a classification problem but work with range sensor data and generate an exhaustive ground truth best view data to train a classifier.
They use a small dataset of 12 objects and CNN based classfier which takes the partial voxels as input to predict the next view. 
In contrast, we use RGB images and do not generate the ground truth but propose a classification based approach which is guided by the supervision from 3D reconstruction directly.
Yang et al. \cite{yang2018active} propose a recurrent network for NBV prediction using passive cameras but use a highly restricted set of viewpoints on the view sphere unlike the broad viewpoints considered in NBV literature and restrained acquisition process.
